% Figure: the worked frame. Horizontal member AC pinned to the wall at A,
% loaded by W at the free end C, and held by a two-force tie BD running
% from a wall anchor D down to the midspan point B of AC. AC is the
% multi-force member; BD is the two-force tie.
\begin{figure}[htbp]
\centering
\begin{tikzpicture}[>=latex, x=1.7cm, y=1.7cm]
  % wall
  \draw[thick] (0,-0.4) -- (0,2.3);
  \foreach \y in {-0.3,0.0,...,2.2} { \draw[enggray, thin] (0,\y) -- (-0.22,\y-0.16); }
  \coordinate (A) at (0,0);    % pin to wall
  \coordinate (B) at (1.5,0);  % tie attaches to AC here
  \coordinate (C) at (3,0);    % free end, load
  \coordinate (D) at (0,2);    % tie anchor on wall
  % multi-force member AC
  \draw[engblue, very thick] (A) -- (C);
  % two-force tie BD
  \draw[engamber, very thick] (D) -- (B);
  % joints
  \fill[engblack] (A) circle (1.6pt); \node[font=\scriptsize, anchor=north east] at (A) {$A$};
  \fill[engblack] (B) circle (1.6pt); \node[font=\scriptsize, anchor=north] at ($(B)+(0,-0.04)$) {$B$};
  \fill[engblack] (C) circle (1.6pt); \node[font=\scriptsize, anchor=north west] at (C) {$C$};
  \fill[engblack] (D) circle (1.6pt); \node[font=\scriptsize, anchor=south east] at (D) {$D$};
  % load W at C
  \draw[->, very thick, engred] (C) -- ++(0,-0.7);
  \node[engred, font=\scriptsize, anchor=west] at ($(C)+(0.05,-0.4)$) {$W=6\,\text{kN}$};
  % tie force direction on AC at B (toward D, tension)
  \draw[->, thick, englime!70!black] (B) -- ($(B)!0.5!(D)$);
  \node[englime!55!black, font=\scriptsize, anchor=west] at ($(B)!0.28!(D)+(0.04,0)$) {$S$};
  % 3-4-5 marks on the tie geometry
  \draw[<->, enggray, thin] (-0.32,0) -- (-0.32,2);
  \node[enggray, font=\scriptsize, anchor=east] at (-0.32,1) {$2\,\si{\meter}$};
  \draw[<->, enggray, thin] (0,-0.55) -- (1.5,-0.55);
  \node[enggray, font=\scriptsize, anchor=north] at (0.75,-0.55) {$1.5\,\si{\meter}$};
  \draw[<->, enggray, thin] (1.5,-0.8) -- (3,-0.8);
  \node[enggray, font=\scriptsize, anchor=north] at (2.25,-0.8) {$1.5\,\si{\meter}$};
\end{tikzpicture}
\caption[Frame with a multi-force member]{The worked frame. The
horizontal member $AC$ is pinned to the wall at $A$, carries the load
$W=6\,\text{kN}$ at its free end $C$, and is held up by the tie $BD$
running from a wall anchor $D$ ($2\,\si{\meter}$ above $A$) down to the
midspan point $B$. Member $AC$ is a multi-force member: it carries three
non-collinear forces (the pin reaction at $A$, the tie force at $B$, and
the load at $C$), so it bends and cannot be treated as a two-force
member. The tie $BD$, loaded only at its two ends, is a two-force member
whose force acts along its own axis. The tie runs a $3$-$4$-$5$ geometry,
$1.5\,\si{\meter}$ horizontal over $2\,\si{\meter}$ vertical, length
$2.5\,\si{\meter}$. Disassembling the frame and taking moments about $A$
isolates the tie force at once, and the tie carries $15\,\text{kN}$,
two and a half times the load it supports.}
\label{fig:vol03ch02:frame-multiforce}
\end{figure}
